\documentclass[10pt]{scrartcl}
\usepackage[mdthm,colorsec,noauthor,noasy]{evan}
\usepackage{handout}
\usepackage{figure-ink}
\newcommand{\aug}[2]{\left[\begin{array}{@{}#1@{}}#2\end{array}\right]}
\newcommand{\vv}[1]{\mathbf{#1}}
\DeclareMathOperator{\spn}{Span}\DeclareMathOperator{\Nul}{Nul}\DeclareMathOperator{\Col}{Col}\DeclareMathOperator{\Row}{Row}\DeclareMathOperator{\tr}{tr}
\title{Math 54 --- Week 8}
\subtitle{Diagonalization; eigenvectors and linear transformations; complex eigenvalues (Lay 5.3--5.5)}
\author{}\date{}\ihead{}
\begin{document}
\maketitle

\section{Facts used from Weeks 1--7}

\begin{enumerate}[label=(W\arabic*),leftmargin=2.8em,itemsep=4pt]
  \ii \target{fact-w1}\textbf{Columns and diagonal scaling (\xweek{math54}{2026-09-14-matrix-algebra-inverses}{Week~4}):} $AP=[\,A\vv p_1\ \cdots\ A\vv p_n\,]$, and for $D=\operatorname{diag}(d_1,\dots,d_n)$, $PD=[\,d_1\vv p_1\ \cdots\ d_n\vv p_n\,]$. The inverse of a $2\times2$ matrix with $ad-bc\ne0$ is $\frac{1}{ad-bc}\begin{psmallmatrix}d&-b\\-c&a\end{psmallmatrix}$.

  \ii \target{fact-w2}\textbf{Invertible matrix theorem \& Rank theorem (\xlink{math54}{2026-09-28-bases-dimension-rank}{thm-lay-4-14}{Weeks~4--6}):} For $n\times n$ matrix $P$: $P$ is invertible $\iff$ its columns are linearly independent $\iff$ its columns form a basis of $\RR^n \iff \det P\ne 0$. $\det(AB)=\det A\det B$. Rank theorem: $\rank A+\dim\Nul A=n$.

  \ii \target{fact-w3}\textbf{Eigenvalues and similarity (\xlink{math54}{2026-10-05-coordinates-eigenvalues}{name-characteristic-equation}{Week~7}):} $\lambda$ is an eigenvalue of $A \iff \det(A-\lambda I)=0$; the eigenspace is $\Nul(A-\lambda I)$. Eigenvectors for distinct eigenvalues are linearly independent (\xlink{math54}{2026-10-05-coordinates-eigenvalues}{thm-lay-5-2}{Lay 5.1 Thm.~2}). Similar matrices ($B=P^{-1}AP$) share the same characteristic polynomial (\xlink{math54}{2026-10-05-coordinates-eigenvalues}{thm-lay-5-4}{Lay 5.2 Thm.~4}). Eigenvalues of a triangular matrix are its diagonal entries (\xlink{math54}{2026-10-05-coordinates-eigenvalues}{thm-lay-5-1}{Lay 5.1 Thm.~1}). $\tr A$ is the diagonal sum.

  \ii \target{fact-w4}\textbf{Coordinates (\xlink{math54}{2026-10-05-coordinates-eigenvalues}{thm-lay-4-9}{Week~7}):} For a basis $\mathcal B$ of $\RR^n$ with $P_{\mathcal B}=[\,\vv b_1\ \cdots\ \vv b_n\,]$: $\vv x=P_{\mathcal B}[\vv x]_{\mathcal B}$ and $[\vv x]_{\mathcal B}=P_{\mathcal B}^{-1}\vv x$. Coordinates preserve linear relations.

  \ii \target{fact-w5}\textbf{Determinants \& dimension (\xlink{math54}{2026-09-28-bases-dimension-rank}{thm-lay-4-9}{Weeks~4--6}):} $\det$ can be expanded along any column. In a $d$-dimensional space, $>d$ vectors are dependent, and any independent set extends to a basis.
\end{enumerate}

\section{This week}

\subsection{Diagonalization (Lay 5.3)}

\begin{qbox}
  The difference equation $\vv x_{k+1}=A\vv x_k$ has closed-form solution $\vv x_k=A^k\vv x_0$. For
  \[
    A=\begin{bmatrix}7&2\\-4&1\end{bmatrix},
  \]
  multiplying $A^{100}$ directly is hopeless. When can $A$ be factored so that $A^k$ is computed algebraically in one step for every $k$?
\end{qbox}

Powers of a diagonal matrix $D=\operatorname{diag}(d_1,\dots,d_n)$ are immediate: $D^k=\operatorname{diag}(d_1^k,\dots,d_n^k)$. Whenever $A=PDP^{-1}$ is \emph{similar} to a diagonal matrix, intermediate $P^{-1}P$ factors cancel in any power:
\[
  A^k = (PDP^{-1})(PDP^{-1})\cdots(PDP^{-1}) = P D^k P^{-1}.
\]

\begin{example*}[Computing $A^k$ for the matrix in the question]
  \begin{steps}
    \ii Eigenvalues: $\tr A=8$ and $\det A=7+8=15$, so
    \[
      \det(A-\lambda I)=\lambda^2-8\lambda+15=(\lambda-5)(\lambda-3),
    \]
    and the eigenvalues are $\lambda=5,3$.
    \ii Eigenvectors: Row~$1$ of $A-5I$ is $(2,2)\implies\vv p_1=(1,-1)$; Row~$1$ of $A-3I$ is $(4,2)\implies\vv p_2=(1,-2)$.
    \ii Setting $P=\begin{psmallmatrix}1&1\\-1&-2\end{psmallmatrix}$ ($\det P=-1$), $D=\begin{psmallmatrix}5&0\\0&3\end{psmallmatrix}$, $P^{-1}=\begin{psmallmatrix}2&1\\-1&-1\end{psmallmatrix}$, we have $AP=PD$, so $A=PDP^{-1}$ and for all $k\ge 0$:
    \[
      A^k = P\begin{bmatrix}5^k&0\\0&3^k\end{bmatrix}P^{-1}
      = \begin{bmatrix}2\cdot 5^k-3^k & 5^k-3^k\\ 2\cdot 3^k-2\cdot 5^k & 2\cdot 3^k-5^k\end{bmatrix}.
    \]
  \end{steps}
\end{example*}

\noindent\textbf{Definition.} A square matrix $A$ is \vocab{diagonalizable} if $A=PDP^{-1}$ for some invertible matrix $P$ and diagonal matrix $D$.

\begin{theorem*}[The Diagonalization Theorem (Lay 5.3, Thm.~5)]\target{thm-lay-5-5}
  An $n\times n$ matrix $A$ is diagonalizable if and only if $A$ has $n$ linearly independent eigenvectors. In fact, $A=PDP^{-1}$ with $D$ diagonal if and only if the columns of $P$ are $n$ linearly independent eigenvectors of $A$, and the diagonal entries of $D$ are the corresponding eigenvalues in the same order.
\end{theorem*}
\begin{ideabox}
  Clear the inverse: $A=PDP^{-1} \iff AP=PD$. Compare columns on both sides using (\ilink{fact-w1}{W1}): column~$j$ of $AP$ is $A\vv p_j$, whereas column~$j$ of $PD$ is $d_j\vv p_j$. Thus $AP=PD$ states that every column $\vv p_j$ satisfies $A\vv p_j=d_j\vv p_j$, while $P$ being invertible is equivalent to its columns $\vv p_1,\dots,\vv p_n$ being linearly independent (and hence nonzero).
\end{ideabox}
\begin{proof}
  Let $P=[\,\vv p_1\ \cdots\ \vv p_n\,]$ and $D=\operatorname{diag}(\lambda_1,\dots,\lambda_n)$. By (\ilink{fact-w1}{W1}), $AP=[\,A\vv p_1\ \cdots\ A\vv p_n\,]$ and $PD=[\,\lambda_1\vv p_1\ \cdots\ \lambda_n\vv p_n\,]$.
  \begin{steps}
    \ii ($\Rightarrow$) If $A=PDP^{-1}$, then $AP=PD$, so $A\vv p_j=\lambda_j\vv p_j$ for each $j$. Since $P$ is invertible, its columns are linearly independent by (\ilink{fact-w2}{W2}) (so none is $\vv 0$), giving $n$ independent eigenvectors.
    \ii ($\Leftarrow$) If $\vv p_1,\dots,\vv p_n$ are independent eigenvectors with $A\vv p_j=\lambda_j\vv p_j$, then $AP=PD$ and $P$ is invertible by (\ilink{fact-w2}{W2}); multiplying on the right by $P^{-1}$ yields $A=PDP^{-1}$.
  \end{steps}
\end{proof}

\begin{fig}
\begin{tikzpicture}[>=Stealth]
  % (a) Before A: parallelogram spanned by eigenvectors p1, p2
  \begin{scope}[xshift=0cm,yshift=0cm,scale=0.78]
    \path[pattern=inkgrain,pattern color=inkblue] (0,0) -- (1,-1) -- (2,-3) -- (1,-2) -- cycle;
    \draw[ink line,figaxis] (-0.6,0) -- (3.0,0) node[right,font=\scriptsize] {\InkLabel{$x_1$}};
    \draw[ink line,figaxis] (0,-3.95) -- (0,1.05) node[above,font=\scriptsize] {\InkLabel{$x_2$}};

    % Subtle dotted eigenlines and parallelogram guides
    \draw[ink guide,draw=black!38] (-0.4,0.4) -- (2.2,-2.2);
    \draw[ink guide,draw=black!38] (-0.3,0.6) -- (1.55,-3.1);
    \draw[ink guide,draw=black!48] (1,-1) -- (2,-3) -- (1,-2);

    % Eigenvectors and focal sum x
    \draw[ink line,figvec,draw=black,line width=0.65pt] (0,0) -- (1,-1) node[above right=2pt,font=\scriptsize] {\InkLabel{$\vv p_1$}};
    \draw[ink line,figvec,draw=black,line width=0.65pt] (0,0) -- (1,-2) node[below left=2pt,font=\scriptsize] {\InkLabel{$\vv p_2$}};
    \draw[ink line,figvec,draw=black,line width=0.65pt] (0,0) -- (2,-3) node[right=4pt,font=\scriptsize] {\InkLabel{$\vv x$}};

    \node[font=\footnotesize] at (1.25,-4.3) {\InkLabel{(a)}};
  \end{scope}

  % Clean center mapping arrow
  \draw[ink line,-{Stealth[length=4.2pt,width=2.4pt]},black,line width=0.65pt] (3.6,-1.1) -- node[above=2pt,font=\footnotesize] {\InkLabel{$A$}} (4.8,-1.1);

  % (b) After A (drawn at 1/3 scale)
  \begin{scope}[xshift=6.05cm,yshift=0cm,scale=0.78]
    \path[pattern=inkgrain,pattern color=inkblue] (0,0) -- ({5/3},{-5/3}) -- ({8/3},{-11/3}) -- (1,-2) -- cycle;
    \draw[ink line,figaxis] (-0.6,0) -- (3.25,0) node[right,font=\scriptsize] {\InkLabel{$x_1$}};
    \draw[ink line,figaxis] (0,-3.95) -- (0,1.05) node[above,font=\scriptsize] {\InkLabel{$x_2$}};

    \draw[ink guide,draw=black!38] (-0.4,0.4) -- (2.3,-2.3);
    \draw[ink guide,draw=black!38] (-0.3,0.6) -- (1.45,-2.9);
    \draw[ink guide,draw=black!48] ({5/3},{-5/3}) -- ({8/3},{-11/3}) -- (1,-2);

    \draw[ink line,figvec,draw=black,line width=0.65pt] (0,0) -- ({5/3},{-5/3}) node[above right=2pt,font=\scriptsize] {\InkLabel{$5\vv p_1$}};
    \draw[ink line,figvec,draw=black,line width=0.65pt] (0,0) -- (1,-2) node[below left=2pt,font=\scriptsize] {\InkLabel{$3\vv p_2$}};
    \draw[ink line,figvec,draw=black,line width=0.65pt] (0,0) -- ({8/3},{-11/3}) node[right=4pt,font=\scriptsize] {\InkLabel{$A\vv x$}};

    \node[font=\footnotesize] at (1.45,-4.3) {\InkLabel{(b)}};
  \end{scope}
\end{tikzpicture}
\figcap{Diagonalization of $A=\begin{psmallmatrix}7&2\\-4&1\end{psmallmatrix}$: \textbf{(a)} $\vv x=\vv p_1+\vv p_2$ in the eigenvector basis $\{\vv p_1,\vv p_2\}$; \textbf{(b)} its image $A\vv x=5\vv p_1+3\vv p_2$ (drawn at $1/3$ scale), stretching $\vv p_1$ by $5$ and $\vv p_2$ by $3$, so $A^k(\vv p_1+\vv p_2)=5^k\vv p_1+3^k\vv p_2$.}
\end{fig}

\method{How to diagonalize a matrix and compute its powers.}
\begin{steps}
  \ii Find the eigenvalues from $\det(A-\lambda I)=0$, with multiplicities.
  \ii For each eigenvalue, find a basis of the eigenspace $\Nul(A-\lambda I)$ (parametric vector form).
  \ii If the bases together contain $n$ vectors, put them as the columns of $P$ and the matching eigenvalues, in the same order, on the diagonal of $D$. If they contain fewer than $n$, $A$ is not diagonalizable (\ilink{thm-lay-5-7}{Thm.~5.7}).
  \ii Check $AP=PD$ column by column ($A\vv p_j=d_j\vv p_j$); this avoids computing $P^{-1}$.
  \ii For powers, $A^k=PD^kP^{-1}$, with $D^k$ obtained by raising each diagonal entry to the $k$-th power.
\end{steps}
\emph{Edge cases.}
\begin{itemize}
  \ii $n$ distinct eigenvalues: diagonalizable without further work (\ilink{thm-lay-5-6}{Thm.~5.6}).
  \ii A repeated eigenvalue of multiplicity $m$: diagonalizable only if its eigenspace has dimension $m$, i.e.\ $\rank(A-\lambda I)=n-m$.
  \ii An eigenvalue $0$ is allowed: $D$ then has a $0$ on the diagonal and $A$ is not invertible, but it may still be diagonalizable.
  \ii Non-real roots of the characteristic polynomial: not diagonalizable with real $P,D$ (Section 2.3 handles the $2\times2$ case).
\end{itemize}

\begin{example*}[A repeated eigenvalue that is diagonalizable]
  Let
  \[
    A = \begin{bmatrix}1&3&3\\-3&-5&-3\\3&3&1\end{bmatrix}.
  \]
  Subtracting row~$3$ from row~$1$ of $A-\lambda I$ factors out $2+\lambda$, giving $\det(A-\lambda I)=-(\lambda-1)(\lambda+2)^2$: eigenvalues $\lambda=1$ and $\lambda=-2$ (multiplicity~$2$). Row reducing $A-I$ gives rank~$2$ and eigenvector $\vv v_1=(1,-1,1)$. For $\lambda=-2$, every row of $A+2I$ is a multiple of $(1,1,1)$, so $\Nul(A+2I)$ is the $2$-plane $x_1+x_2+x_3=0$ with basis $\vv v_2=(-1,1,0)$, $\vv v_3=(-1,0,1)$, giving three vectors in total. Thus $A=PDP^{-1}$ with
  \[
    P = \begin{bmatrix}1&-1&-1\\-1&1&0\\1&0&1\end{bmatrix}, \qquad
    D = \begin{bmatrix}1&0&0\\0&-2&0\\0&0&-2\end{bmatrix}, \qquad (\det P=1\ne 0).
  \]
\end{example*}

\begin{theorem*}[Distinct eigenvalues guarantee diagonalizability (Lay 5.3, Thm.~6)]\target{thm-lay-5-6}
  An $n\times n$ matrix with $n$ distinct eigenvalues is diagonalizable.
\end{theorem*}
\begin{proof}
  Choose one eigenvector for each distinct eigenvalue; by \xlink{math54}{2026-10-05-coordinates-eigenvalues}{thm-lay-5-2}{Lay 5.1 Thm.~2} (\ilink{fact-w3}{W3}) these $n$ vectors are linearly independent, so \ilink{thm-lay-5-5}{Thm.~5.5} applies. (The condition is sufficient, not necessary: $I$ and the example \emph{A repeated eigenvalue that is diagonalizable} are diagonalizable with repeated eigenvalues.)
\end{proof}

When an eigenvalue $\lambda$ has algebraic multiplicity $m>1$, the only obstacle to diagonalizability is an eigenspace of dimension strictly less than $m$.

\begin{theorem*}[Eigenspace dimension vs.\ multiplicity (Lay 5.3, Thm.~7a)]\target{thm-lay-5-7a}
  If $\lambda_0$ is a real eigenvalue of $A$ with algebraic multiplicity $m$, then $1\le \dim\Nul(A-\lambda_0 I)\le m$.
\end{theorem*}
\begin{ideabox}
  Extend a basis $\vv v_1,\dots,\vv v_p$ of $E(\lambda_0,A)$ to a basis of $\RR^n$ forming columns of $P$. Since $P^{-1}AP$ sends the first $p$ standard basis vectors to $\lambda_0\vv e_j$, $\det(P^{-1}AP-\lambda I)$ factors out $(\lambda_0-\lambda)^p$ by cofactor expansion down the first $p$ columns, forcing $(\lambda-\lambda_0)^p \mid \det(A-\lambda I)$.
\end{ideabox}
\begin{proof}
  Let $\vv v_1,\dots,\vv v_p$ ($p\ge 1$) be a basis of $\Nul(A-\lambda_0 I)$, extend to a basis $\vv v_1,\dots,\vv v_n$ of $\RR^n$ (\ilink{fact-w5}{W5}), set $P=[\,\vv v_1\ \cdots\ \vv v_n\,]$, and let $B=P^{-1}AP$. For $j\le p$, column~$j$ of $B$ is $P^{-1}A\vv v_j=\lambda_0\vv e_j$. Expanding $\det(B-\lambda I)$ repeatedly down its first $p$ columns gives $\det(B-\lambda I)=(\lambda_0-\lambda)^p q(\lambda)$. By similarity (\ilink{fact-w3}{W3}), $\det(A-\lambda I)=\det(B-\lambda I)$, so $\lambda_0$ has multiplicity $m\ge p$.
\end{proof}

\begin{theorem*}[Counting eigenvectors (Lay 5.3, Thm.~7)]\target{thm-lay-5-7}\target{thm-lay-5-7b}
  Let $A$ be $n\times n$ with distinct real eigenvalues $\lambda_1,\dots,\lambda_r$ of multiplicities $m_1,\dots,m_r$, and let $\mathcal B_k$ be a basis of $\Nul(A-\lambda_k I)$ of dimension $d_k$.
  \begin{enumerate}
    \ii[(a)] The combined collection $\mathcal B_1\cup\dots\cup\mathcal B_r$ is linearly independent.
    \ii[(b)] $A$ is diagonalizable $\iff \sum_{k=1}^r d_k = n \iff \det(A-\lambda I)$ factors completely into real linear factors AND $d_k=m_k$ for every $k$.
  \end{enumerate}
\end{theorem*}
\begin{proof}
  \begin{steps}
    \ii \textbf{(a).} Suppose a linear combination of $\bigcup\mathcal B_k$ equals $\vv 0$. Group the terms from $\mathcal B_k$ into $\vv w_k\in\Nul(A-\lambda_k I)$, so $\vv w_1+\dots+\vv w_r=\vv 0$. Any nonzero $\vv w_k$ would be eigenvectors for distinct eigenvalues summing to $\vv 0$, impossible by (\ilink{fact-w3}{W3}); thus each $\vv w_k=\vv 0$, forcing all weights in each basis $\mathcal B_k$ to be $0$.
    \ii \textbf{(b).} By (a) and (\ilink{fact-w5}{W5}), the maximum number of independent eigenvectors in any eigenspace selection is $\sum_{k=1}^r d_k$, so \ilink{thm-lay-5-5}{Thm.~5.5} applies $\iff \sum d_k=n$. Since $\sum d_k \le \sum m_k \le \deg\det(A-\lambda I)=n$ (\ilink{thm-lay-5-7a}{Thm.~5.7a}; the multiplicities of distinct roots add up to at most the degree), equality holds $\iff \sum m_k=n$ (complete real splitting) and $d_k=m_k$ for every $k$.
  \end{steps}
\end{proof}

\begin{example*}[Same eigenvalues, not diagonalizable]
The matrix
    \[
      A = \begin{bmatrix}2&4&3\\-4&-6&-3\\3&3&1\end{bmatrix}
    \]
    also has $\det(A-\lambda I)=-(\lambda-1)(\lambda+2)^2$, yet $A+2I$ has rank~$2$, so $\dim\Nul(A+2I)=3-2=1<2$. Total independent eigenvectors $=1+1=2<3$, so $A$ is \emph{not} diagonalizable.
\end{example*}

\begin{example*}[A triangular $4\times4$ matrix that is diagonalizable]
The matrix
    \[
      A = \begin{bmatrix}2&0&0&0\\0&2&0&0\\1&3&4&0\\0&0&0&4\end{bmatrix}
    \]
    has eigenvalues $2,2,4,4$. Here $\rank(A-2I)=2 \implies \dim E(2)=2$ (basis $(-3,1,0,0),(-2,0,1,0)$) and $\rank(A-4I)=2 \implies \dim E(4)=2$ (basis $\vv e_3,\vv e_4$). Since $2+2=4$, $A$ is diagonalizable.
\end{example*}

\revisit
$A$ can be factored as $A=PDP^{-1}$ exactly when it has $n$ independent eigenvectors (\ilink{thm-lay-5-5}{Thm.~5.5}); then $A^k=PD^kP^{-1}$ for every $k$. For the matrix in the question, the eigenvalues $5,3$ are distinct, and the example \emph{Computing $A^k$ for the matrix in the question} gives $A^{100}$ by putting $k=100$ in a fixed formula: for instance its $(1,1)$ entry is $2\cdot5^{100}-3^{100}$.

\begin{shortcuts}
  \ii \textbf{Checking $A=PDP^{-1}$:} Check $AP=PD$ ($A\vv p_j=d_j\vv p_j$) instead of computing $P^{-1}$.
  \ii \textbf{Diagonalizability test via rank:} Simple eigenvalues always contribute $1$; a repeated eigenvalue $\lambda$ of multiplicity $m$ works iff $\rank(A-\lambda I)=n-m$ (Rank Thm.).
  \ii \textbf{Limits $A^k$ as $k\to\infty$:} If every eigenvalue satisfies $\lambda=1$ or $|\lambda|<1$, then $D^k\to\operatorname{diag}(1,\dots,0)$ gives $\lim_{k\to\infty}A^k$.
\end{shortcuts}

\begin{trapbox}
  \ii \textbf{Column order of $P$ and $D$ must match:} If column~$1$ of $P$ is a $5$-eigenvector, $D_{11}$ must be $5$.
  \ii \textbf{Diagonalizable $\neq$ invertible:} $\operatorname{diag}(1,0)$ is diagonalizable and singular; $\begin{psmallmatrix}1&1\\0&1\end{psmallmatrix}$ is invertible and \emph{not} diagonalizable.
  \ii \textbf{Repeated eigenvalues can go either way:} $2I$ is diagonal; $\begin{psmallmatrix}2&1\\0&2\end{psmallmatrix}$ is not.
\end{trapbox}

\subsection{Eigenvectors and linear transformations (Lay 5.4)}

\begin{qbox}
  In the example \emph{Computing $A^k$ for the matrix in the question}, $A=PDP^{-1}$ with $D=\operatorname{diag}(5,3)$. What does this factorization reveal about the geometric transformation $\vv x\mapsto A\vv x$ itself when viewed in the basis $\mathcal B=\{(1,-1),(1,-2)\}$? And how do we represent linear maps between abstract vector spaces like $\mathbb P_2$ by matrices?
\end{qbox}

Let $T\colon V\to W$ be linear, with basis $\mathcal B=\{\vv b_1,\dots,\vv b_n\}$ of $V$ and basis $\mathcal C$ of $W$. Expanding $\vv x=\sum_{j=1}^n r_j\vv b_j$ and using linearity of $T$ and $[\cdot]_{\mathcal C}$ (\ilink{fact-w4}{W4}),
\[
  [T(\vv x)]_{\mathcal C} = r_1[T(\vv b_1)]_{\mathcal C}+\dots+r_n[T(\vv b_n)]_{\mathcal C} = M[\vv x]_{\mathcal B},
  \qquad M = \bigl[\,[T(\vv b_1)]_{\mathcal C}\ \ \cdots\ \ [T(\vv b_n)]_{\mathcal C}\,\bigr].
\]
When $W=V$ and $\mathcal C=\mathcal B$, $M=[T]_{\mathcal B}$ is the \vocab{$\mathcal B$-matrix of $T$}, satisfying $[T(\vv x)]_{\mathcal B}=[T]_{\mathcal B}[\vv x]_{\mathcal B}$.

\begin{fig}
\begin{tikzpicture}[>=Stealth]
  \node[font=\small] (V) at (0,1.95) {\InkLabel{$\vv x\in V$}};
  \node[font=\small] (W) at (4.6,1.95) {\InkLabel{$T(\vv x)\in W$}};
  \node[font=\small] (Rn) at (0,0) {\InkLabel{$[\vv x]_{\mathcal B}\in\RR^n$}};
  \node[font=\small] (Rm) at (4.6,0) {\InkLabel{$[T(\vv x)]_{\mathcal C}\in\RR^m$}};

  \draw[ink line,-{Stealth[length=4.0pt,width=2.3pt]},black,line width=0.6pt,shorten <=6pt,shorten >=6pt] (V) -- node[above=2pt,font=\footnotesize] {\InkLabel{$T$}} (W);
  \draw[ink line,-{Stealth[length=3.8pt,width=2.2pt]},black!80,line width=0.5pt,shorten <=5pt,shorten >=5pt] (V) -- node[left=4pt,font=\scriptsize] {\InkLabel{$[\cdot]_{\mathcal B}$}} (Rn);
  \draw[ink line,-{Stealth[length=3.8pt,width=2.2pt]},black!80,line width=0.5pt,shorten <=5pt,shorten >=5pt] (W) -- node[right=4pt,font=\scriptsize] {\InkLabel{$[\cdot]_{\mathcal C}$}} (Rm);
  \draw[ink line,-{Stealth[length=4.0pt,width=2.3pt]},black,line width=0.6pt,shorten <=6pt,shorten >=6pt] (Rn) -- node[below=2pt,font=\footnotesize] {\InkLabel{$M$}} (Rm);
\end{tikzpicture}
\figcap{The matrix $M$ of a linear transformation $T\colon V\to W$ relative to bases $\mathcal B$ and $\mathcal C$ satisfies $[T(\vv x)]_{\mathcal C} = M[\vv x]_{\mathcal B}$, so the diagram commutes.}
\end{fig}

\begin{example*}[Differentiation on $\mathbb P_2$]
  Let $D\colon\mathbb P_2\to\mathbb P_2$ be $D(p)=p'$ in $\mathcal B=\{1,t,t^2\}$. Since $D(1)=0$, $D(t)=1$, $D(t^2)=2t$,
  \[
    [D]_{\mathcal B} = \begin{bmatrix}0&1&0\\0&0&2\\0&0&0\end{bmatrix}, \qquad
    [D(5+3t-t^2)]_{\mathcal B} = [D]_{\mathcal B}\begin{bmatrix}5\\3\\-1\end{bmatrix} = \begin{bmatrix}3\\-2\\0\end{bmatrix}\leftrightarrow 3-2t.
  \]
  Because coordinates keep all linear relations (\ilink{fact-w4}{W4}), $\vv x\neq\vv 0$ satisfies $T(\vv x)=\lambda\vv x \iff [T]_{\mathcal B}[\vv x]_{\mathcal B}=\lambda[\vv x]_{\mathcal B}$. Here $[D]_{\mathcal B}$ is upper-triangular with diagonal $0,0,0$, so $0$ is the only eigenvalue (constant polynomials).
\end{example*}

\method{How to find the matrix of a linear transformation relative to bases.}
\begin{steps}
  \ii Apply $T$ to each vector $\vv b_j$ of the basis $\mathcal B$ of the domain.
  \ii Write each $T(\vv b_j)$ in coordinates relative to the basis $\mathcal C$ of the codomain.
  \ii Use $[T(\vv b_j)]_{\mathcal C}$ as column $j$. The matrix is $m\times n$ for $\dim V=n$, $\dim W=m$.
  \ii For $T(\vv x)=A\vv x$ on $\RR^n$ with $\mathcal C=\mathcal B$ given by the columns of $P$, the $\mathcal B$-matrix is $P^{-1}AP$ (\ilink{thm-lay-5-8}{Thm.~5.8} below).
\end{steps}
\emph{Edge cases.} The number of columns is the dimension of the domain, not of the codomain. To make the $\mathcal B$-matrix diagonal, choose $\mathcal B$ to consist of eigenvectors.

\begin{example*}[The matrix of a map from polynomials to $\RR^2$]
$T\colon\mathbb P_2\to\RR^2$, $T(p)=(p(0),\ p(1))$, with $\mathcal B=\{1,t,t^2\}$ and the standard basis of $\RR^2$: $T(1)=(1,1)$, $T(t)=(0,1)$, $T(t^2)=(0,1)$, so
\[
  M=\begin{bmatrix}1&0&0\\1&1&1\end{bmatrix}.
\]
For $p=2-t+3t^2$: $M(2,-1,3)=(2,\ 2-1+3)=(2,4)$, and indeed $p(0)=2$, $p(1)=4$.
\end{example*}

\begin{theorem*}[Diagonal and $\mathcal B$-matrix representation (Lay 5.4, Thm.~8)]\target{thm-lay-5-8}
  Suppose $A=PCP^{-1}$ with $P$ invertible, and let $\mathcal B$ be the basis of $\RR^n$ formed by the columns of $P$. Then $C$ is the $\mathcal B$-matrix of the transformation $\vv x\mapsto A\vv x$. In particular, if $A=PDP^{-1}$ with $D$ diagonal, then $D$ is the $\mathcal B$-matrix of $\vv x\mapsto A\vv x$.
\end{theorem*}
\begin{proof}
  By (\ilink{fact-w4}{W4}), $[A\vv x]_{\mathcal B}=P^{-1}(A\vv x) = P^{-1}AP[\vv x]_{\mathcal B} = C[\vv x]_{\mathcal B}$ for every $\vv x\in\RR^n$. Applying both $C$ and $[\vv x\mapsto A\vv x]_{\mathcal B}$ to $[\vv b_j]_{\mathcal B}=\vv e_j$ confirms their columns match.
\end{proof}

\revisit In the eigenvector basis $\mathcal B=\{(1,-1),(1,-2)\}$, the $\mathcal B$-matrix of $\vv x\mapsto A\vv x$ is $D=\operatorname{diag}(5,3)$ (\ilink{thm-lay-5-8}{Thm.~5.8}): the transformation scales coordinate~$1$ by $5$ and coordinate~$2$ by $3$ independently. A linear map between abstract spaces such as $\mathbb P_2$ is represented by the matrix whose columns are the coordinate vectors $[T(\vv b_j)]_{\mathcal C}$, as for differentiation above.

\begin{example*}[A non-diagonalizable $2\times2$ matrix in an almost-diagonal basis]
  $A=\begin{psmallmatrix}4&-9\\4&-8\end{psmallmatrix}$ has $\det(A-\lambda I)=(\lambda+2)^2$ and $E(-2,A)=\spn\{(3,2)\}$ ($\rank(A+2I)=1$). Using $\mathcal B=\{(3,2),(2,1)\}$,
  \[
    [x\mapsto A\vv x]_{\mathcal B} = P^{-1}AP = \begin{bmatrix}-1&2\\2&-3\end{bmatrix}\begin{bmatrix}4&-9\\4&-8\end{bmatrix}\begin{bmatrix}3&2\\2&1\end{bmatrix} = \begin{bmatrix}-2&1\\0&-2\end{bmatrix},
  \]
  upper-triangular with $1$ on the superdiagonal (a $2\times 2$ Jordan block).
\end{example*}

\begin{proposition*}[Facts about similar matrices]\target{prop-similar-facts}
  Let $B=P^{-1}AP$. Then $\det B=\det A$; $\tr B=\tr A$; $B^k=P^{-1}A^kP$; $A=(P^{-1})^{-1}BP^{-1}$ (similarity is symmetric), and if also $C=Q^{-1}BQ$ then $C=(PQ)^{-1}A(PQ)$ (transitive); and if $A\vv v=\lambda\vv v$, then $B(P^{-1}\vv v)=\lambda(P^{-1}\vv v)$.
\end{proposition*}
\begin{proof}
  \begin{steps}
    \ii $\det B=\det P^{-1}\det A\det P=\det A$, since $\det P^{-1}\det P=\det I=1$ (\ilink{fact-w2}{W2}).
    \ii \textbf{$\tr(XY)=\tr(YX)$} for $n\times n$ $X,Y$: both equal $\sum_i\sum_k x_{ik}y_{ki}$. Then $\tr B=\tr\bigl((P^{-1}A)P\bigr)=\tr\bigl(P(P^{-1}A)\bigr)=\tr A$.
    \ii In $B^k=(P^{-1}AP)(P^{-1}AP)\cdots(P^{-1}AP)$ each inner $PP^{-1}$ is $I$, leaving $P^{-1}A^kP$.
    \ii Multiplying $B=P^{-1}AP$ by $P$ on the left and $P^{-1}$ on the right gives $A=PBP^{-1}$; and $C=Q^{-1}P^{-1}APQ=(PQ)^{-1}A(PQ)$.
    \ii $B(P^{-1}\vv v)=P^{-1}AP P^{-1}\vv v=P^{-1}A\vv v=\lambda P^{-1}\vv v$, and $P^{-1}\vv v\ne\vv 0$ because $P^{-1}$ is invertible.
  \end{steps}
\end{proof}

\begin{shortcuts}
  \ii \textbf{$\mathcal B$-matrix on $\RR^n$:} $[T]_{\mathcal B}=P^{-1}AP$ where columns of $P$ are $\mathcal B$.
  \ii \textbf{Eigenvalues of an abstract $T\colon V\to V$:} Compute eigenvalues of $[T]_{\mathcal B}$ in any convenient basis $\mathcal B$.
\end{shortcuts}

\begin{trapbox}
  \ii \textbf{Order is $P^{-1}AP$, \emph{not} $PAP^{-1}$:} $PAP^{-1}=\begin{psmallmatrix}3&0\\6&5\end{psmallmatrix}\neq D$ in the example \emph{Computing $A^k$ for the matrix in the question}.
  \ii \textbf{Similar matrices have different eigenvectors in general:} $D$ has eigenvector $\vv e_1$, while $A$ has $\vv p_1=P\vv e_1$.
\end{trapbox}

\subsection{Complex eigenvalues (Lay 5.5)}

\begin{qbox}
  The $90^\circ$ counterclockwise rotation $A=\begin{psmallmatrix}0&-1\\1&0\end{psmallmatrix}$ moves every nonzero real vector off its own line, so $\det(A-\lambda I)=\lambda^2+1=0$ has no real roots. What geometric action do its \emph{complex} eigenvalues $\pm i$ encode?
\end{qbox}

Allowing scalars in $\CC$, $\lambda\in\CC$ is a \vocab{complex eigenvalue} of $A$ if $A\vv x=\lambda\vv x$ for some nonzero $\vv x\in\CC^n$. The definitions and proofs of Weeks~1--7 use only addition, subtraction, multiplication and division of scalars, which obey the same rules for complex numbers, so they hold over $\CC$ without change; in particular $\lambda$ is a complex eigenvalue exactly when $\det(A-\lambda I)=0$.

\begin{theorem*}[Fundamental theorem of algebra]\target{thm-fta}
  Every polynomial of degree $n\ge1$ with complex coefficients factors as $c(z-z_1)\cdots(z-z_n)$ with $c,z_1,\dots,z_n\in\CC$. Consequently every $n\times n$ matrix has $n$ complex eigenvalues counted with multiplicity.
\end{theorem*}
\begin{proof}
  The proof is in \ilink{proof-fta}{Proof of the fundamental theorem of algebra} at the end. The second sentence follows by applying the first to the characteristic polynomial, which has degree $n$ (\xlink{math54}{2026-10-05-coordinates-eigenvalues}{prop-charpoly-degree}{The characteristic polynomial has degree $n$}).
\end{proof}

Taking complex conjugates preserves sums and products ($\overline{z+w}=\bar z+\bar w$ and $\overline{zw}=\bar z\,\bar w$, so entry by entry $\overline{B\vv x}=\overline{B}\,\overline{\vv x}$); for a \emph{real} matrix $A=\overline{A}$,
\[
  A\vv x=\lambda\vv x \implies A\overline{\vv x}=\overline{\lambda}\,\overline{\vv x},
\]
so complex eigenvalues of a real matrix come in \emph{conjugate pairs} $\lambda,\overline{\lambda}$ with conjugate eigenvectors. For $A=\begin{psmallmatrix}0&-1\\1&0\end{psmallmatrix}$, $\lambda=i$ has eigenvector $\vv x=(1,-i)$, and automatically $\lambda=-i$ has $\overline{\vv x}=(1,i)$.

\medskip\noindent\textbf{Rotation--scaling matrices.} For real $a,b$ (not both $0$), define
\[
  C = \begin{bmatrix}a&-b\\b&a\end{bmatrix}
  = r\begin{bmatrix}\cos\varphi&-\sin\varphi\\\sin\varphi&\cos\varphi\end{bmatrix},
  \qquad r = |\lambda|=\sqrt{a^2+b^2}=\sqrt{\det C},\quad \cos\varphi=\frac{a}{r},\ \sin\varphi=\frac{b}{r}.
\]
Its eigenvalues are $\lambda=a\pm bi$, because $\det(C-\lambda I)=(a-\lambda)^2+b^2=0$ means $a-\lambda=\pm bi$. The second form of $C$ is the rotation matrix of Week~3 multiplied by $r$, so $C$ rotates by $\varphi$ and scales lengths by $r$. The modulus $r$ is the \textbf{scale factor}, and the argument $\varphi$ of $(a,b)$ is the \textbf{rotation angle}. (Watch the quadrant: $\begin{psmallmatrix}1&1\\-1&1\end{psmallmatrix}$ has $a=1,b=-1$, so $(1,-1)$ in Quadrant~IV gives $\varphi=-45^\circ$, whereas $\begin{psmallmatrix}\sqrt{3}&-1\\1&\sqrt{3}\end{psmallmatrix}$ gives $r=2,\varphi=30^\circ$.)

\begin{fig}
\begin{tikzpicture}[>=Stealth]
  % (a) Complex plane: lambda = a + bi = r e^{i phi}
  \begin{scope}[xshift=0cm,yshift=0cm,scale=0.92]
    \draw[ink line,figaxis] (-0.45,0) -- (3.0,0) node[right,font=\scriptsize] {\InkLabel{$\operatorname{Re}$}};
    \draw[ink line,figaxis] (0,-0.45) -- (0,2.15) node[above,font=\scriptsize] {\InkLabel{$\operatorname{Im}$}};
    \draw[ink guide,draw=black!45] (1.9,0) -- (1.9,1.2) -- (0,1.2);
    \draw[ink line,figvec,draw=black,line width=0.65pt,shorten >=2.2pt] (0,0) -- (1.9,1.2) node[midway,sloped,above=2pt,font=\scriptsize] {\InkLabel{$r$}};
    \node[figdot] at (1.9,1.2) {};
    \node[above right=2pt,font=\scriptsize] at (1.9,1.2) {\InkLabel{$a+bi$}};
    \node[below=2pt,font=\scriptsize] at (1.9,0) {\InkLabel{$a$}};
    \node[left=2pt,font=\scriptsize] at (0,1.2) {\InkLabel{$b$}};
    \draw[ink line,black!75,line width=0.45pt] (0.64,0) arc (0:32.3:0.64);
    \node[font=\scriptsize] at (0.88,0.24) {\InkLabel{$\varphi$}};
    \node[font=\footnotesize] at (1.25,-0.75) {\InkLabel{(a)}};
  \end{scope}

  % (b) Geometric action of C in R^2
  \begin{scope}[xshift=6.0cm,yshift=0cm,scale=0.92]
    \draw[ink line,figaxis] (-0.45,0) -- (3.0,0) node[right,font=\scriptsize] {\InkLabel{$x_1$}};
    \draw[ink line,figaxis] (0,-0.45) -- (0,2.15) node[above,font=\scriptsize] {\InkLabel{$x_2$}};
    \draw[ink guide,draw=black!40] (1,0) arc (0:32.3:1);
    \draw[ink line,figvec,draw=black,line width=0.65pt] (0,0) -- (1,0) node[below=2pt,font=\scriptsize] {\InkLabel{$\vv e_1$}};
    \draw[ink line,figvec,draw=black,line width=0.65pt] (0,0) -- (1.9,1.2) node[pos=0.62,sloped,above=2pt,font=\scriptsize] {\InkLabel{$r$}};
    \node[above right=2pt,font=\scriptsize] at (1.9,1.2) {\InkLabel{$C\vv e_1=(a,b)$}};
    \draw[ink line,-{Stealth[length=3.4pt,width=2.0pt]},black!75,line width=0.45pt] (0.64,0) arc (0:32.3:0.64);
    \node[font=\scriptsize] at (1.3,0.3) {\InkLabel{$\varphi$}};
    \node[font=\footnotesize] at (1.25,-0.75) {\InkLabel{(b)}};
  \end{scope}
\end{tikzpicture}
\figcap{\textbf{(a)} Complex eigenvalue $\lambda=a+bi=r(\cos\varphi+i\sin\varphi)$ with modulus $r=\sqrt{a^2+b^2}$ and argument $\varphi$. \textbf{(b)} The matrix $C=\begin{psmallmatrix}a&-b\\b&a\end{psmallmatrix}$ rotates every vector by $\varphi$ and scales length by $r$, sending $\vv e_1$ to $(a,b)$.}
\end{fig}

\revisit The $90^\circ$ rotation matrix has $a=0,b=1$: its eigenvalues $\pm i$ have modulus $r=1$ (no scaling) and argument $\varphi=90^\circ$ (the rotation angle).

\begin{theorem*}[Real $2\times2$ matrices with complex eigenvalues (Lay 5.5, Thm.~9)]\target{thm-lay-5-9}
  Let $A$ be a real $2\times2$ matrix with a complex eigenvalue $\lambda=a-bi$ ($b\ne0$) and associated eigenvector $\vv v\in\CC^2$. Then
  \[
    A = P C P^{-1}, \qquad\text{where}\qquad P = [\,\operatorname{Re}\vv v\ \ \operatorname{Im}\vv v\,],\qquad C = \begin{bmatrix}a&-b\\b&a\end{bmatrix}.
  \]
\end{theorem*}
\begin{ideabox}
  Write $\vv v=\vv x+i\vv y$ with $\vv x,\vv y\in\RR^2$. Since $A$ is real, splitting $A(\vv x+i\vv y)=(a-bi)(\vv x+i\vv y)$ into real and imaginary parts gives two real equations: $A\vv x = a\vv x + b\vv y$ and $A\vv y = -b\vv x + a\vv y$. By (\ilink{fact-w1}{W1}), those two columns together state $AP = PC$. (Choosing $\lambda=a-bi$ instead of $a+bi$ flips $\operatorname{Im}\vv v \mapsto -\operatorname{Im}\vv v$ and ensures $+b$ sits in entry $(2,1)$ of $C$.)
\end{ideabox}
\begin{proof}
  Let $\vv x=\operatorname{Re}\vv v$ and $\vv y=\operatorname{Im}\vv v$. Since $A$ is real, $A\vv v=A\vv x+iA\vv y$, while
  \[
    \lambda\vv v=(a-bi)(\vv x+i\vv y)=(a\vv x+b\vv y)+i(-b\vv x+a\vv y).
  \] Equating real and imaginary parts gives $A\vv x=a\vv x+b\vv y$ and $A\vv y=-b\vv x+a\vv y$, i.e.\ $AP=PC$ (column by column, \ilink{fact-w1}{W1}).
  \begin{steps}
    \ii \textbf{$\vv x$ and $\vv y$ are independent.} Suppose not. If $\vv x=\vv 0$, then $\vv v=i\vv y$ with $\vv y$ real and nonzero, and $A\vv v=\lambda\vv v$ gives $A\vv y=\lambda\vv y$. If $\vv x\ne\vv 0$ and $\vv y=c\vv x$ with $c$ real, then $\vv v=(1+ic)\vv x$ and dividing $A\vv v=\lambda\vv v$ by $1+ic\ne0$ gives $A\vv x=\lambda\vv x$. Either way a real nonzero vector $\vv u$ satisfies $A\vv u=\lambda\vv u$. The left side is real; the right side has imaginary part $-b\vv u\ne\vv 0$, since $b\ne0$. This is impossible.
    \ii So $P$ is invertible (\ilink{fact-w2}{W2}), and multiplying $AP=PC$ on the right by $P^{-1}$ gives $A=PCP^{-1}$.
  \end{steps}
\end{proof}

\method{How to find complex eigenvalues, an eigenvector, and $A=PCP^{-1}$ for a real $2\times2$ matrix.}
\begin{steps}
  \ii Solve $\lambda^2-(\tr A)\lambda+\det A=0$; if the discriminant is negative, $\lambda=a\pm bi$ with $a=\tfrac12\tr A$.
  \ii Take $\lambda=a-bi$. A nonzero row $(p\ \ q)$ of $A-\lambda I$ gives the eigenvector $\vv v=(q,-p)$: the row times $\vv v$ is $pq-qp=0$, and the other row is a complex multiple of this one because $A-\lambda I$ is singular and nonzero.
  \ii Set $P=[\,\operatorname{Re}\vv v\ \ \operatorname{Im}\vv v\,]$ and $C=\begin{psmallmatrix}a&-b\\b&a\end{psmallmatrix}$ (\ilink{thm-lay-5-9}{Thm.~5.9}); check $AP=PC$.
  \ii The scale factor is $r=\sqrt{a^2+b^2}=\sqrt{\det A}$, and the angle $\varphi$ satisfies $\cos\varphi=a/r$, $\sin\varphi=b/r$.
\end{steps}
\emph{Edge cases.} Using $\lambda=a+bi$ instead puts $-b$ below the diagonal of $C$; pair the sign choices consistently. Any complex multiple of $\vv v$ is also an eigenvector and gives a different, equally valid $P$. If $b<0$ in $C$, the rotation is clockwise.

\begin{example*}[A rotation--scaling factorization with $r=\sqrt5$]
For $A=\begin{psmallmatrix}1&-2\\1&3\end{psmallmatrix}$, $\tr A=4$ and $\det A=5$, so $\lambda^2-4\lambda+5=0$ gives $\lambda=2\pm i$. For $\lambda=2-i$, row~1 of $A-\lambda I$ is $(-1+i,\ -2)$, so $\vv v=(-2,\ 1-i)$. Then
\[
  P=\begin{bmatrix}-2&0\\1&-1\end{bmatrix},\qquad C=\begin{bmatrix}2&-1\\1&2\end{bmatrix},\qquad
  AP=\begin{bmatrix}-4&2\\1&-3\end{bmatrix}=PC.
\]
The scale factor is $r=\sqrt5$ and $\varphi$ is the angle of $(2,1)$, about $26.6^\circ$.
\end{example*}

\begin{example*}[An elliptical orbit (Lay 5.5, Ex.~2)]
  For
  \[
    A = \begin{bmatrix}0.5&-0.6\\0.75&1.1\end{bmatrix},
  \]
  $\tr A=1.6$, $\det A=1$, so $\lambda^2-1.6\lambda+1=0 \implies \lambda=0.8\pm 0.6i$ ($|\lambda|=\sqrt{0.64+0.36}=1$). Using $\lambda=0.8-0.6i$: row~$1$ of $A-\lambda I$ is $(-0.3+0.6i,\ -0.6)$, giving eigenvector $\vv v=(2,-1+2i)$. Taking $\operatorname{Re}\vv v=(2,-1)$ and $\operatorname{Im}\vv v=(0,2)$ gives
  \[
    P=\begin{bmatrix}2&0\\-1&2\end{bmatrix},\qquad
    C=\begin{bmatrix}0.8&-0.6\\0.6&0.8\end{bmatrix},\qquad
    A=PCP^{-1}.
  \]
  Since $r=1$, $C$ is pure rotation by $\varphi=\arccos(0.8)\approx 36.9^\circ$. Iterating $\vv x_k=A^k\vv x_0 = P C^k\vv e_1$ (from $\vv x_0=\operatorname{Re}\vv v=P\vv e_1$) traces an \textbf{ellipse} $\{\cos\theta(2,-1)+\sin\theta(0,2)\}$ advancing $\approx 36.9^\circ$ per step. (If $r>1$ it spirals outward; if $r<1$ it spirals inward.) For $3\times3$ matrices like $\operatorname{diag}(C,1.07)$, trajectories spiral while expanding/shrinking along the third eigenvector.
\end{example*}

\begin{fig}
\begin{tikzpicture}[>=Stealth]
  % (a) Pure rotation C on the unit circle in B-coordinates
  \begin{scope}[xshift=0cm,yshift=0cm,scale=1.22]
    \path[pattern=inkgrain,pattern color=inkblue] (0,0) circle (1.0);
    \draw[ink line,figaxis] (-1.45,0) -- (1.65,0) node[right,font=\scriptsize] {\InkLabel{$u_1$}};
    \draw[ink line,figaxis] (0,-1.45) -- (0,1.55) node[above,font=\scriptsize] {\InkLabel{$u_2$}};
    \InkStroke[black,line width=0.65pt]{(0,0) circle (1.0)}

    % Basis vectors e1, e2 (crisp black ink clearly distinct from hairline axes)
    \InkStroke[figvec,draw=black,line width=0.65pt,shorten >=2.6pt]{(0,0) -- (1,0)}
    \node[font=\scriptsize] at (0.55,-0.19) {\InkLabel{$\vv e_1$}};
    \InkStroke[figvec,draw=black,line width=0.65pt,shorten >=2.6pt]{(0,0) -- (0,1)}
    \node[font=\scriptsize] at (-0.21,0.52) {\InkLabel{$\vv e_2$}};

    % Rotation angle phi = 36.87 deg
    \draw[ink guide,draw=black!45] (0,0) -- (0.8,0.6);
    \draw[ink line,-{Stealth[length=3.4pt,width=2.0pt]},black!75,line width=0.45pt] (0.46,0) arc (0:36.87:0.46);
    \node[font=\scriptsize] at (0.65,0.16) {\InkLabel{$\varphi$}};

    % Iterates C^k e1
    \foreach \k/\lbl/\pos in {0/0/below right, 1/1/right, 2/2/above right, 3/3/above left} {
      \pgfmathsetmacro{\ux}{cos(\k*36.87)}
      \pgfmathsetmacro{\uy}{sin(\k*36.87)}
      \node[figdot] at (\ux,\uy) {};
      \node[\pos=2.5pt,font=\scriptsize] at (\ux,\uy) {\InkLabel{$[\vv x_{\lbl}]_{\mathcal B}$}};
    }
    \node[font=\footnotesize] at (0.05,-1.72) {\InkLabel{(a)}};
  \end{scope}

  % Connecting similarity map P
  \draw[ink line,-{Stealth[length=4.2pt,width=2.4pt]},black,line width=0.65pt] (2.65,0) -- node[above=2pt,font=\scriptsize] {\InkLabel{$P=[\,\operatorname{Re}\vv v\ \ \operatorname{Im}\vv v\,]$}} (4.85,0);

  % (b) Invariant ellipse in R^2
  \begin{scope}[xshift=7.55cm,yshift=0cm,scale=0.76]
    \path[pattern=inkgrain,pattern color=inkblue,domain=0:360,samples=100,variable=\t] plot ({2*cos(\t)},{-cos(\t)+2*sin(\t)}) -- cycle;
    \draw[ink line,figaxis] (-2.6,0) -- (2.9,0) node[right,font=\scriptsize] {\InkLabel{$x_1$}};
    \draw[ink line,figaxis] (0,-2.5) -- (0,2.75) node[above,font=\scriptsize] {\InkLabel{$x_2$}};
    \InkStroke[black,line width=0.65pt,domain=0:360,samples=100,variable=\t]{plot ({2*cos(\t)},{-cos(\t)+2*sin(\t)})}

    % Conjugate semi-axes Re v = (2,-1) and Im v = (0,2) (consistent solid peer vectors)
    \InkStroke[figvec,draw=black,line width=0.65pt,shorten >=3.2pt]{(0,0) -- (2,-1)}
    \node[font=\scriptsize] at (0.85,-0.98) {\InkLabel{$\operatorname{Re}\vv v$}};
    \InkStroke[figvec,draw=black,line width=0.65pt,shorten >=3.2pt]{(0,0) -- (0,2)}
    \node[left=3pt,font=\scriptsize] at (0,1.05) {\InkLabel{$\operatorname{Im}\vv v$}};

    % Iterates x0..x3 cleanly outside ellipse
    \foreach \k/\lbl/\pos in {0/0/below right, 1/1/right, 2/2/above right, 3/3/above left} {
      \pgfmathsetmacro{\px}{2*cos(\k*36.87)}
      \pgfmathsetmacro{\py}{-cos(\k*36.87)+2*sin(\k*36.87)}
      \node[figdot] at (\px,\py) {};
      \node[\pos=3pt,font=\scriptsize] at (\px,\py) {\InkLabel{$\vv x_{\lbl}$}};
    }
    \node[font=\footnotesize] at (0.1,-2.76) {\InkLabel{(b)}};
  \end{scope}
\end{tikzpicture}
\figcap{Orbit $\vv x_k=A^k\vv x_0$ starting at $\vv x_0=\operatorname{Re}\vv v=(2,-1)$: \textbf{(a)} in the $C$-basis, $C^k\vv e_1$ steps around the circle by $\varphi\approx 36.9^\circ$; \textbf{(b)} $P=[\,\operatorname{Re}\vv v\ \ \operatorname{Im}\vv v\,]$ maps that circle to the invariant ellipse (spiraling in when $r<1$, out when $r>1$).}
\end{fig}

\begin{shortcuts}
  \ii \textbf{Complex eigenvector of $2\times2$:} Any nonzero row $(p\ \; q)$ of $A-\lambda I$ immediately yields eigenvector $(q,-p)$.
  \ii \textbf{Rotation--scaling parameters:} From $\lambda=a\pm bi$, $r=\sqrt{a^2+b^2}=\sqrt{\det A}$ and $\varphi$ is the angle of $(a,b)$ (mind the quadrant).
\end{shortcuts}

\begin{trapbox}
  \ii \textbf{Columns of $P$ in \ilink{thm-lay-5-9}{Thm.~5.9} are $\operatorname{Re}\vv v$ and $\operatorname{Im}\vv v$ (real vectors),} \emph{not} $\vv v$ and $\overline{\vv v}$, paired with $\lambda=a-bi$.
  \ii \textbf{Check $\sin\varphi=b/r$ as well as $\cos\varphi=a/r$:} $\cos\varphi$ alone leaves $\pm\varphi$ ambiguous.
  \ii \textbf{Only \emph{real} matrices force conjugate eigenvalue pairs:} A complex matrix like $\operatorname{diag}(i,2)$ need not have $-i$ as an eigenvalue.
\end{trapbox}

\section{Homework this week}
HW7, Sections 5.3--5.4, and HW8, Section 5.5.
\begin{itemize}
  \ii \textbf{Section 5.3:} \#1, \#3 ($A^k=PD^kP^{-1}$); \#7 ($2\times 2$ distinct eigenvalues $\implies$ \ilink{thm-lay-5-6}{Thm.~5.6}); \#13, \#19 (check $\rank(A-\lambda I)$ for repeated $\lambda$, \ilink{thm-lay-5-7}{Thm.~5.7}); \#22, \#23, \#24, \#26 (\ilink{thm-lay-5-5}{Thm.~5.5}, the example \emph{Same eigenvalues, not diagonalizable}, $\operatorname{diag}(1,0)$, $I$); \#30 ($1+1<3$, \ilink{thm-lay-5-7}{Thm.~5.7}).
  \ii \textbf{Section 5.4:} \#1, \#3, \#4 (columns are $[T(\vv b_j)]_{\mathcal C}$); \#5 ($[T(\vv x)]_{\mathcal B}=[T]_{\mathcal B}[\vv x]_{\mathcal B}$); \#7, \#8, \#9 ($P^{-1}AP$); \#15, \#20 ($T(p)=\lambda p$); \#22, \#23, \#27 (similarity facts).
  \ii \textbf{Section 5.5:} \#3, \#5 ($2\times 2$ complex eigenvalues/eigenvectors); \#10 ($r$ and $\varphi$, watch quadrant); \#13, \#15 (\ilink{thm-lay-5-9}{Thm.~5.9} with $\lambda=a-bi$); \#24 (ellipse in Lay 5.5 Ex.~2).
\end{itemize}

\section{Problem set}

\begin{prob}\target{prob-P1}
  \textbf{P1.}\ Diagonalize the matrix below, and find a formula for $A^k$.
  \[
    A = \begin{bmatrix}2&3\\4&1\end{bmatrix}.
  \]
\end{prob}

\begin{prob}\target{prob-P2}
  \textbf{P2.}\ Decide whether each matrix is diagonalizable, and diagonalize it if it is:
  \[
    A=\begin{bmatrix}2&0&0\\1&2&0\\0&0&3\end{bmatrix},\qquad
    B=\begin{bmatrix}4&0&-2\\2&5&4\\0&0&5\end{bmatrix}.
  \]
\end{prob}

\begin{prob}\target{prob-P3}
  \textbf{P3.}\ A $4\times 4$ matrix has eigenvalues $1,1,2,3$, and $\rank(A-I)=3$. Is $A$ diagonalizable? What if $\rank(A-I)=2$?
\end{prob}

\begin{prob}\target{prob-P4}
  \textbf{P4.}\ Let $T\colon\mathbb P_1\to\mathbb P_2$ be $T(p)=(t+1)p(t)$. Find the matrix of $T$ relative to $\mathcal B=\{1,t\}$ and $\mathcal C=\{1,t,t^2\}$, and use it to compute $T(3-2t)$.
\end{prob}

\begin{prob}\target{prob-P5}
  \textbf{P5.}\ For the matrix $A$ below and $\mathcal B=\{(1,1),(1,2)\}$, find the $\mathcal B$-matrix of $\vv x\mapsto A\vv x$. What does it tell you about $(1,1)$ and $(1,2)$?
  \[
    A = \begin{bmatrix}3&-1\\2&0\end{bmatrix}.
  \]
\end{prob}

\begin{prob}\target{prob-P6}
  \textbf{P6.}\ For the matrix below, find its complex eigenvalues, an eigenvector for $3-i$, and $P,C$ with $A=PCP^{-1}$. What are the scale factor and the rotation angle of $C$?
  \[
    A = \begin{bmatrix}2&-1\\2&4\end{bmatrix}.
  \]
\end{prob}

\begin{prob}\target{prob-P7}
  \textbf{P7.}\ Prove: if $A$ is diagonalizable and invertible, then $A^{-1}$ is diagonalizable. If $A$ is diagonalizable, is $A^{\mathsf T}$?
\end{prob}

\begin{prob}\target{prob-P8}
  \textbf{P8.}\ For the matrix below and $\vv x_0=(1,0)$, find $\vv x_k=A^k\vv x_0$ and its limit as $k\to\infty$.
  \[
    A = \begin{bmatrix}3/4&1/4\\1/4&3/4\end{bmatrix}.
  \]
\end{prob}

\section{Proofs of the facts used}

\subsection{Proof of the fundamental theorem of algebra}\target{proof-fta}
This proves \ilink{thm-fta}{the fundamental theorem of algebra}. We use two facts from calculus: a polynomial $p(z)$, viewed as a function of the two real variables $x,y$ in $z=x+iy$, is continuous, and a continuous real function on a closed disk attains a minimum value there.

\begin{steps}
  \ii \textbf{Polar form and $k$-th roots.} Every nonzero $w\in\CC$ is $w=\rho(\cos\theta+i\sin\theta)$ with $\rho=|w|>0$. Multiplying two such numbers multiplies the moduli and adds the angles (expand and use the angle-addition formulas). Hence $\bigl(\rho^{1/k}(\cos\tfrac{\theta}{k}+i\sin\tfrac{\theta}{k})\bigr)^k=w$: every nonzero complex number has a $k$-th root.
  \ii \textbf{$|p(z)|$ is large far from $0$.} Let $p(z)=a_nz^n+\dots+a_0$ with $a_n\ne0$, $n\ge1$. Let $S=|a_{n-1}|+\dots+|a_0|$. For $|z|\ge1$ each $|z|^j\le|z|^{n-1}$ ($j<n$), so
  \[
    |a_{n-1}z^{n-1}+\dots+a_0|\le S|z|^{n-1},\qquad |p(z)|\ge|a_n||z|^n-S|z|^{n-1}=|z|^{n-1}(|a_n||z|-S).
  \] Choose $R\ge1$ with $|a_n|R-S\ge|a_0|+1$; then $|p(z)|>|a_0|=|p(0)|$ whenever $|z|\ge R$.
  \ii \textbf{A minimum exists.} On the closed disk $|z|\le R$, $|p|$ attains a minimum at some $z_0$. Since $0$ is in the disk, $|p(z_0)|\le|p(0)|<|p(z)|$ for $|z|\ge R$, so $|p(z_0)|\le|p(z)|$ for every $z\in\CC$.
  \ii \textbf{The minimum value is $0$.} Suppose $p(z_0)\ne0$. Expanding $p(z_0+w)$ in powers of $w$ gives $p(z_0+w)=p(z_0)+cw^k+w^{k+1}g(w)$, where $c\ne0$ is the first nonzero coefficient after the constant term ($k\ge1$ exists because $p$ is not constant) and $g$ is a polynomial. By Step~1 choose $u$ with $u^k=-p(z_0)/c$, and set $w=tu$ with $0<t<1$. Then
  \[
    p(z_0+tu)=p(z_0)(1-t^k)+t^{k+1}u^{k+1}g(tu),\qquad |p(z_0+tu)|\le|p(z_0)|(1-t^k)+t^{k+1}M,
  \]
  where $M$ bounds $|u^{k+1}g(tu)|$ for $0\le t\le1$. For $0<t<1$ with $tM<|p(z_0)|$ the right side is less than $|p(z_0)|$ (subtract $|p(z_0)|$ and divide by $t^k$), contradicting Step~3. So $p(z_0)=0$.
  \ii \textbf{Factor theorem.} If $p(z_0)=0$, then $p(z)=p(z)-p(z_0)=\sum_j a_j(z^j-z_0^j)$, and $z^j-z_0^j=(z-z_0)(z^{j-1}+z^{j-2}z_0+\dots+z_0^{j-1})$, so $p(z)=(z-z_0)q(z)$ with $\deg q=n-1$ and the same leading coefficient $a_n$.
  \ii \textbf{Induction.} Applying Steps~3--5 to $q$, then to its quotient, and so on, after $n$ steps $p(z)=a_n(z-z_1)\cdots(z-z_n)$.
\end{steps}

\end{document}
